\section{Năm nhóm tín hiệu hiệu quả của SourceVerify}
Hình \ref{fig:method-groups} chia hệ thống phân tích ảnh thành 5 nhóm tín hiệu chính: nguồn gốc tệp, miền tần số, nhiễu/kết cấu, biên cạnh/thống kê và màu sắc/nén ảnh.

\begin{figure}[h]
\centering
\begin{tikzpicture}[
  group/.style={rectangle,draw=Navy,rounded corners,fill=LightBlue,align=center,minimum width=3.6cm,minimum height=1.0cm},
  item/.style={rectangle,draw=Blue,rounded corners,fill=white,align=center,minimum width=3.4cm,minimum height=0.65cm}
]
\node[group] at (0,0) {Nguồn gốc tệp};
\node[item] at (0,-1.0) {Metadata Analysis};
\node[group] at (4,0) {Miền tần số};
\node[item] at (4,-1.0) {Spectral Nyquist};
\node[item] at (4,-1.8) {Multi-scale Reconstruction};
\node[group] at (8,0) {Nhiễu / kết cấu};
\node[item] at (8,-1.0) {Noise Residual};
\node[item] at (8,-1.8) {Gradient Micro-Texture};
\node[group] at (2,-3.2) {Biên cạnh / thống kê};
\node[item] at (2,-4.2) {Edge Coherence};
\node[item] at (2,-5.0) {Benford's Law};
\node[group] at (6.5,-3.2) {Màu sắc / nén ảnh};
\node[item] at (6.5,-4.2) {Chromatic Aberration};
\node[item] at (6.5,-5.0) {DCT Block Artifacts};
\node[item] at (6.5,-5.8) {Color Channel Correlation};
\end{tikzpicture}
\caption{Phân nhóm 5 tín hiệu ảnh hiệu quả trong báo cáo}
\label{fig:method-groups}
\end{figure}

\subsection{Metadata Analysis}
Metadata là thông tin mô tả đi kèm tệp ảnh, ví dụ tên phần mềm, thời gian tạo, kích thước, thiết bị chụp, thông tin EXIF hoặc profile màu. Nếu metadata ghi nhận phần mềm tạo ảnh như một công cụ AI, đây là tín hiệu mạnh. Ngược lại, nếu ảnh có thông tin camera hợp lý, điều đó có thể ủng hộ giả thuyết ảnh thật.

Tuy nhiên, metadata rất dễ bị xoá hoặc sửa. Do đó, SourceVerify không dùng metadata làm bằng chứng duy nhất. Method này phù hợp để phát hiện các trường hợp rõ ràng, đồng thời hỗ trợ người dùng hiểu nguồn gốc tệp ban đầu.

\subsection{Spectral Nyquist Analysis}
Phân tích Nyquist tập trung vào miền tần số của ảnh. Ảnh sinh bởi AI hoặc ảnh được upscale có thể để lại các mẫu lặp ở tần số cao, nhất là gần biên Nyquist. Những mẫu này không phải lúc nào cũng thấy bằng mắt thường nhưng có thể xuất hiện khi phân tích gradient hoặc phổ tần số.

\subsection{Multi-scale Reconstruction}
Multi-scale Reconstruction đánh giá sự ổn định của ảnh ở nhiều mức phân giải. Ảnh tự nhiên thường có biến thiên cục bộ phức tạp; khi giảm và khôi phục kích thước, sai số tái tạo có thể thay đổi theo vùng. Ảnh AI đôi khi có sai số đồng đều hơn do cấu trúc được sinh ra từ mô hình.

\subsection{Noise Residual}
Noise Residual tách phần nhiễu còn lại sau khi loại bỏ thành phần mượt của ảnh. Camera thật thường tạo nhiễu cảm biến không hoàn toàn đồng nhất. Nhiễu còn phụ thuộc ánh sáng, ISO, cảm biến và pipeline xử lý ảnh. Ảnh AI có thể quá sạch hoặc có nhiễu không giống nhiễu vật lý.

\subsection{Edge Coherence}
Edge Coherence phân tích các đường biên và vùng chuyển tiếp. Ảnh thật thường có biên cạnh chịu ảnh hưởng bởi vật thể, ánh sáng, chuyển động, lens blur và sharpening. Ảnh AI có thể tạo đường biên quá mượt, thiếu vi sai hoặc không nhất quán giữa các vùng.

\subsection{Gradient Micro-Texture}
Gradient Micro-Texture đánh giá các thay đổi rất nhỏ trong vùng chuyển sắc và vùng bề mặt. Ảnh thật thường có micro-texture do cảm biến, chất liệu vật thể, ánh sáng và nhiễu tự nhiên. Ảnh AI, đặc biệt ảnh được làm mịn, có thể thiếu loại kết cấu nhỏ này.

\subsection{Benford's Law}
Benford's Law là quy luật thống kê về phân bố chữ số đầu trong nhiều dữ liệu tự nhiên. Trong ảnh, ý tưởng tương tự có thể áp dụng cho phân bố độ lớn gradient hoặc các đặc trưng số học khác. Nếu phân bố lệch quá nhiều so với mẫu tự nhiên, ảnh có thể đã qua tổng hợp hoặc xử lý bất thường.

\subsection{Chromatic Aberration}
Chromatic Aberration là sai lệch màu nhỏ do ống kính thật gây ra, thường xuất hiện ở vùng biên có độ tương phản cao. Camera thật và ống kính vật lý có thể để lại viền màu nhẹ giữa các kênh RGB. Ảnh AI không nhất thiết có dấu vết quang học này.

\subsection{DCT Block Artifacts}
JPEG sử dụng biến đổi DCT theo khối. Quá trình nén tạo ra artifact đặc trưng, nhất là ở ảnh đã lưu nhiều lần hoặc có chất lượng nén thấp. Ảnh thật thường có dấu vết nén phù hợp với lịch sử lưu ảnh, trong khi ảnh AI có thể có vết nén quá đều, thiếu tự nhiên hoặc không tương thích với metadata.

\subsection{Color Channel Correlation}
Trong ảnh tự nhiên, ba kênh màu R, G, B có tương quan do ánh sáng, bề mặt vật thể, cảm biến và pipeline xử lý ảnh. Ảnh AI có thể sinh kênh màu theo cách không hoàn toàn giống camera thật, dẫn đến tương quan bất thường hoặc quá đều.

\begin{longtable}{p{0.06\textwidth}p{0.22\textwidth}p{0.30\textwidth}p{0.24\textwidth}p{0.15\textwidth}}
\caption{Tóm tắt 5 nhóm tín hiệu hiệu quả của SourceVerify}\label{tab:signal-groups}\\
\toprule
\textbf{STT} & \textbf{Nhóm tín hiệu} & \textbf{Method đại diện} & \textbf{Dấu hiệu phát hiện} & \textbf{Hạn chế chính} \\
\midrule
\endfirsthead
\toprule
\textbf{STT} & \textbf{Nhóm tín hiệu} & \textbf{Method đại diện} & \textbf{Dấu hiệu phát hiện} & \textbf{Hạn chế chính} \\
\midrule
\endhead
1 & Nguồn gốc tệp & Metadata Analysis & Phần mềm tạo ảnh, EXIF, kích thước, tên tệp & Dễ bị xoá/sửa \\
2 & Miền tần số & Spectral Nyquist, Multi-scale Reconstruction & Đỉnh tần số, artifact upsampling, sai số tái tạo bất thường & Bị ảnh hưởng bởi resize \\
3 & Nhiễu/kết cấu & Noise Residual, Gradient Micro-Texture & Nhiễu quá sạch, thiếu micro-texture tự nhiên & Nhạy với nén/filter \\
4 & Biên cạnh/thống kê & Edge Coherence, Benford's Law & Biên quá mượt, phân bố gradient lệch tự nhiên & Không kết luận đơn lẻ \\
5 & Màu sắc/nén ảnh & Chromatic Aberration, DCT Block Artifacts, Color Channel Correlation & Sai lệch màu, vết nén JPEG, tương quan RGB bất thường & Phụ thuộc chỉnh sửa ảnh \\
\bottomrule
\end{longtable}

